% \siyin{rewrite!!}
\subsection{Foundation models in robotics}
Recent advancements in End-to-End Vision-Language-Action (VLA) models adopt a co-training paradigm that integrates a pre-trained VLM backbone with a dedicated action expert~\cite{Black20240AV, Intelligence202505AV}. By harnessing the VLM's extensive world knowledge~\cite{Bai2025Qwen3VLTR} to interpret complex scenes, this approach effectively transforms the backbone into a generalist agent capable of mapping raw observations directly to continuous actions. 
Alternatively, the hierarchical paradigm decouples reasoning from execution~\cite{Shi2025HiRO, Li2025HAMSTERHA, Shentu2024FromLT}. In this framework, a high-level policy, typically a VLM, processes visual observations and task descriptions to generate intermediate goals, such as natural language sub-tasks or latent embeddings. These intermediate representations then condition a low-level policy to output continuous actions. Building upon this hierarchical structure, our work introduces a novel hierarchical memory architecture designed to empower the high-level VLM to generate more precise and context-aware instructions for the low-level controller.



%However, despite their success in short-term manipulation, this end-to-end coupling struggles with long-horizon tasks. Constrained by limited context windows and implicit state tracking, the model's ability to maintain a consistent chain of thought degrades as the interaction history grows.
%While hierarchical frameworks attempt to mitigate this by decoupling high-level planning from low-level execution, a critical misalignment remains: Robotic tasks are inherently stateful and temporal, yet in long-horizon scenarios, the VLM planner often fails to track task progress, leading to hallucinations or redundant actions.
%Different from previous approaches that rely solely on the VLM's implicit context, we argue that explicit memory is crucial for long-horizon planning. In this work, we augment the VLM planner with a dedicated memory module to provide persistent state tracking. This design effectively bridges the gap between the VLM's reasoning capabilities and the temporal consistency required for robust, multi-stage robotic tasks.


\subsection{Robotic memory}
Memory capabilities are fundamental for VLA to successfully execute complex, long-horizon manipulation tasks. 
Consequently, extensive research has been dedicated to endowing vla with historical context, though most existing implementations remain predominantly image-based~\cite{Huang2022VisualLM}. 
These works typically follow two paradigms: utilizing external memory banks~\cite{Li2025MAPVLAMP, Torne2025LearningLD, Fang2025SAM2ActIV, Huang2023OutOS, Zhu2020VisionDialogNB, Zheng2024TraceVLAVT, Hu20253DLLMMemLS}, that employ FIFO queues~\cite{Sridhar2025MemERSU} or fused buffers~\cite{Shi2025MemoryVLAPM} to store raw observations, or exploring history compression to distill information into compact representations. Notably, MemER~\cite{Sridhar2025MemERSU} recently advanced the latter by proposing keyframe selection algorithms to better preserve critical information within subtasks.
Distinct from these methods, we propose a cross-modal, explicitly managed memory bank designed to enhance task consistency and effectively retain long-term temporal context.

% Vision-Language-Action (VLA) models have shown promise in robotic manipulation, yet processing long-horizon tasks remains a bottleneck due to the quadratic complexity of attention mechanisms. While simply extending context windows allows for retaining recent history, it imposes prohibitive computational costs. Consequently, methods like fixed-size buffers or token compression are employed to manage input length.
% Notably, MEMER introduces a hierarchical framework that employs keyframe algorithms to compress information within subtasks.
% However, these compression strategies often operate on low-level visual similarity or predefined heuristics, risking the loss of semantic details that are visually subtle but causally pivotal for decision-making.
% Departing from passive history compression, we propose an active memory retrieval mechanism for VLA agents. By querying historical data driven by task uncertainty, our model selectively recalls relevant interactions. This strategy avoids maintaining a large active context, ensuring computational efficiency without compromising long-term reasoning








% However, such time-based management is often rigid and task-agnostic. Rigid segmentation can fragment continuous subtasks or miss sparse but critical events that occur off-boundary, as the memory controller ignores the agent's current goal state. Furthermore, these systems largely treat memory maintenance as a passive logging process rather than an active decision.
% In contrast, we propose a task-driven memory management paradigm. Instead of arbitrary time windows, our approach dynamically structures memory around goal-relevant episodes (e.g., subtask completion). This allows the agent to actively curate its memory—deciding when to aggregate observations and explicitly updating unreliable records—ensuring that memory maintenance is aligned with the agent's planning and execution needs.